\scanpage{63}
% archive foliation: 62; reused typescript, upside down in scan.

\unsure{1.|7.}

- considération la demande des besoins locaux.

4 - \uline{Comment faire de la recherche dans un établissement d'enseignement supérieur}

4.1. \uline{Les établissements et leurs moyens de base} :

Un établissement d'enseignement supérieur possède du personnel enseignant dont une partie seulement de l'activité est consacrée à l'enseignement ; le reste peut donc être consacré à la recherche. Il est difficile d'avancer une estimation chiffrée de ce partage, qui varie selon les disciplines. Considérons en première approximation que le temps consacré à l'activité de recherche est de 40\% en moyenne du temps de l'enseignant chercheur. Par ailleurs, l'établissement dispose de locaux construits et entretenus à l'aide de crédits inscrits sur des chapitres budgétaires de la D.E.S.U.R. (non inclus dans l'enveloppe recherche) Enfin, l'établissement dispose de personnel administratif, technique, ouvrier et de service (A.T.O.S.) dont une part de l'activité est consacrée à la recherche. Il apparait donc d'emblée que chaque établissement possède un potentiel ``hors enveloppe-recherche'', qu'il doit mobiliser pour le consacrer à la recherche.

Pour chiffrer ce budget, hors enveloppe-recherche, qui est pourtant consacré à la recherche, on pourrait prendre 40\% de la masse salariale des enseignants et admettre, en première approximation, que 30\% du personnel A.T.O.S. effectue un travail de recherche. En admettant en outre qu'une part du budget général du secrétariat d'Etat (chapitre 36-11) contribue à la recherche (ne serait-ce que par le paiement des dépenses de chauffage des bâtiments ...), on arrive en 1975 à une approximation des dépenses pour la recherche de l'ordre de 1 100 à 1 200 millions de francs. Il peut donc être constaté que par le simple fait qu'elle rétribue à plein temps du personnel qui fait de l'enseignement à temps partiel, l'ensemble des établissements de la D.E.S.U.R. apporte au budget ``recherche'' une somme supérieure au milliard de francs.

4.2. \uline{Les crédits de la D.E.S.U.R. pour constituer ces moyens}

La D.E.S.U.R. dispose de deux chapitres budgétaires,d'un de ``fonctionnement'', l'autre ``d'équipement'', qui permettent d'aider cet effort de recherche dans chaque établissement. Le tableau ci-après indique l'évolution de ces budgets (millions de francs courants)

